\subsection{Qwen and Llama KV-Cache Comparison}

\subsubsection{Experimental setup}

We evaluated Qwen3.5-9B and Meta-Llama-3.1-8B-Instruct on an
NVIDIA GeForce RTX 5090 using a 65,536-token context window, full GPU
offloading, flash attention, and Hermes Agent. Both configurations used the
same Q8\_0 model weights. The controlled variable within each model pair was
the KV-cache format: the traditional configuration used Q8\_0 for both the
key and value caches, whereas the TurboQuant configuration used Turbo4 for
both caches. Consequently, these experiments evaluate Turbo4 KV-cache
compression rather than TurboQuant weight quantization.

\subsubsection{GPU-memory results}

\begin{table*}[t]
\centering
\small
\caption{Idle GPU-memory allocation with traditional and Turbo4 KV caches.}
\label{tab:combined-kv-memory}
\begin{tabular}{lrrrr}
\toprule
Model & Traditional & Turbo4 & Saving & Reduction \\
\midrule
Qwen3.5-9B & 10,310 MiB & 9,662 MiB & 648 MiB & 6.29\% \\
Llama 3.1 8B & 12,900 MiB & 10,672 MiB & 2,228 MiB & 17.27\% \\
\bottomrule
\end{tabular}
\end{table*}

Turbo4 reduced idle GPU-memory allocation for both models. Qwen3.5-9B
saved 648 MiB, corresponding to a 6.29\% reduction, while Llama 3.1 8B
saved 2,228 MiB, corresponding to a 17.27\% reduction. The larger reduction
for Llama indicates that the benefit of KV-cache compression depends on the
model architecture and cache dimensions. This is the clearest improvement
demonstrated by the experiments.

\subsubsection{Runtime, utilization, and power}

For the Qwen travel-planning task, the traditional configuration completed
in 2,199.93 seconds, whereas Turbo4 required 2,269.75 seconds. Turbo4 was
therefore approximately 69.82 seconds, or 3.17\%, slower in this run.
Average GPU utilization was similar at 92.50\% and 93.33\%, respectively.
Average power also remained effectively unchanged, decreasing from
563.19 W to 561.48 W. Because the Turbo4 run lasted longer, its estimated
energy consumption increased slightly from approximately 0.344 kWh to
0.353 kWh.

For Llama, both configurations repeatedly reached 99--100\% GPU utilization
and consumed approximately 575 W during active inference. However, the
available records did not contain sufficiently consistent execution-time,
average-power, or energy measurements for a controlled Llama speed
comparison. Therefore, no speed improvement can be claimed for Llama.

The Qwen document-analysis run also does not demonstrate a raw inference
speed advantage. The traditional run required approximately 884 seconds,
while the recorded Turbo4 run required 2,504.23 seconds and encountered
output truncation and four context-compression events. This difference
includes agent orchestration, continuation, file operations, and context
management, so it must not be interpreted as a direct measurement of token
generation speed.

\subsubsection{Agent-task quality}

\begin{table*}[t]
\centering
\small
\caption{Summary of the autonomous travel-planning results.}
\label{tab:combined-agent-quality}
\begin{tabular}{llll}
\toprule
Configuration & File generation & Main correctness problem & Result \\
\midrule
Qwen traditional & 9/9 files & Final cost was \texteuro{}9,290 & Failed \\
Qwen Turbo4 & 9/9 files & Detailed cost was about \texteuro{}6,947.50 & Failed \\
Llama traditional & 1/9 files & Paris-only itinerary and unsupported budget & Failed \\
Llama Turbo4 & Claimed complete & Minimum stated cost was \texteuro{}5,800 & Failed \\
\bottomrule
\end{tabular}
\end{table*}

Neither cache configuration consistently satisfied the autonomous travel
task. Both Qwen runs created all nine required files, but both exceeded the
\texteuro{}5,000 budget and contained unreliable final verification. The
traditional Qwen run calculated a total of \texteuro{}9,290. The Turbo4
Qwen run reported \texteuro{}3,968.50, although its detailed CSV values
totaled approximately \texteuro{}6,947.50.

The traditional Llama run performed substantially worse: it created only
one of nine required files and produced a 30-day itinerary limited to Paris,
despite requiring eight European countries. The Llama Turbo4 response was
more complete, but its reported accommodation and transportation costs
already totaled \texteuro{}5,800, exceeding the complete task budget before
food, attractions, and other expenses were included.

The apparent quality differences cannot be attributed confidently to
KV-cache format. Quantization changes memory representation, while agent
quality is also affected by sampling variation, generated output length,
tool-call behavior, conversation state, and context compression.

\subsubsection{Context behavior}

Both model families experienced pressure on the 65,536-token context limit.
Hermes used response continuation and context compression when the
conversation and tool history became too large. Turbo4 reduces the physical
GPU memory occupied by the KV cache, but it does not increase the logical
context limit configured by the server. Consequently, Turbo4 did not prevent
output truncation or context compression.

Repeated context compression can also reduce agent reliability because exact
requirements, file paths, calculations, and completed steps may be shortened
or omitted from the retained summary. This behavior likely contributed to
some of the inconsistencies observed in the generated outputs.

\subsubsection{Limitations}

The GPU-memory comparison is reasonably controlled because each model pair
used the same weights, hardware, context size, and server settings, with the
KV-cache format as the intended changed variable. The end-to-end runtime and
quality comparisons are less conclusive because most configurations were
tested only once, generated-token counts were not matched, tool-call paths
differed, context-compression events were inconsistent, and outputs were not
verified equally.

Furthermore, \texttt{nvidia-smi} measures total GPU allocation rather than
the exact amount of KV cache occupied by active tokens. Time to first token,
prompt-processing throughput, decode throughput, and exact KV-cache buffer
sizes were not measured consistently across all configurations.

\subsubsection{Conclusion}

Turbo4 provided a clear GPU-memory benefit for both evaluated models. It
reduced idle allocation by 6.29\% for Qwen3.5-9B and 17.27\% for Llama 3.1
8B, with Llama saving more than 2.2 GiB. However, the current experiments do
not demonstrate a consistent improvement in end-to-end execution time,
energy consumption, or autonomous-task quality. The evidence therefore
supports Turbo4 primarily as a VRAM-efficiency optimization. Repeated trials
with matched token counts, fresh agent sessions, exact KV-cache measurements,
time-to-first-token measurements, and decode-throughput measurements are
required for a definitive performance comparison.
